\documentclass[10pt]{beamer}

\usepackage[utf8]{inputenc}
\usepackage{amsmath, amssymb, bm}
\usepackage{booktabs}
\usepackage{tikz}
\usetikzlibrary{positioning,arrows.meta,shapes.geometric}

\usetheme{Madrid}
\usecolortheme{default}
\setbeamertemplate{navigation symbols}{}
\setbeamertemplate{footline}[frame number]

\title[Two-Period Model]{Two-Period Model:\\Zoning and Transportation Investment}
\author{Ke Jiang}
\institute{Southern Methodist University}
\date{}

\begin{document}

\begin{frame}
  \titlepage
\end{frame}

\begin{frame}{Purpose of the Model}
\begin{block}{Central mechanism}
Local zoning restrictions limit the housing and density response to improved accessibility. A transportation agency that takes zoning as given may therefore find transit investment less attractive than it would under coordinated zoning and transportation policy.
\end{block}

\vspace{0.3cm}
\begin{itemize}
    \item There are two periods in the model.
    \item In the first period, zoning is chosen locally by municipalities who cares about incumbent interests.
    \item Then transportation investment is chosen at the metropolitan level taking zoning as given.
    \item In the second period,Households sort across residence, workplace.
    \item The economy is closed: welfare changes come from the spatial allocation of residents, jobs, housing prices, wages, congestion, and accessibility.
\end{itemize}
\end{frame}

\begin{frame}{Basic Environment}
\begin{itemize}
    \item A metropolitan area is partitioned into municipalities $g \in G$.
    \item Municipality $g$ contains locations $S^g$:
    \[
        S = \bigcup_{g \in G} S^g,
        \qquad
        S^g \cap S^{g'} = \emptyset \quad \text{for } g \neq g'.
    \]
    \item Each location $i$ has a zoning capacity $h_{0,i}$.
    \item Higher $h_{0,i}$ means looser zoning and greater housing capacity.
    \item Transportation investment $I_{jk}$ on link $(j,k)$ reduces commuting costs and improves accessibility.
\end{itemize}

\vspace{0.25cm}
\begin{block}{Policy tension}
Zoning is local and politically fragmented, while transportation returns are regional and depend on the entire spatial pattern of density.
\end{block}
\end{frame}

\begin{frame}{Key Agents}
\begin{columns}[T]
\column{0.33\textwidth}
\textbf{Municipal landowners}
\begin{itemize}
    \item Choose local housing capacity $h_{0,i}$
    \item Represent incumbent interests
    \item Take other municipalities' zoning as given
\end{itemize}

\column{0.33\textwidth}
\textbf{Transportation agency}
\begin{itemize}
    \item Observes full zoning vector $h_0$
    \item Chooses network investment $I$
    \item Maximizes regional household welfare subject to a budget
\end{itemize}

\column{0.33\textwidth}
\textbf{Households}
\begin{itemize}
    \item Begin at origin $o$
    \item Choose residence $i$
    \item Choose workplace $j$
    \item Face housing prices, wages, moving costs, commuting costs, and congestion
\end{itemize}
\end{columns}
\end{frame}

\begin{frame}{Timing}
\begin{columns}[T]
\column{0.48\textwidth}
\begin{block}{Time $t=0$}
\begin{enumerate}
    \item Each municipality $g$ chooses $\{h_{0,i}\}_{i \in S^g}$.
    \item Municipalities play a Nash game in zoning.
    \item The transportation agency observes $h_0$.
    \item The agency chooses optimal network investment $I$.
\end{enumerate}
\end{block}

\column{0.48\textwidth}
\begin{block}{Time $t=1$}
\begin{enumerate}
    \item Households choose residence $i$, workplace $j$, moving frctions involved
    \item Housing markets clear and prices $q_i$ are determined.
    \item Labor markets clear and wages $w_j$ are determined.
    \item Local populations $N_i$ are determined.
\end{enumerate}
\end{block}
\end{columns}

\vspace{0.25cm}\end{frame}

\begin{frame}{Municipal Zoning Choice}
A representative landowner in municipality $g$ chooses local housing capacity:
\[
u_{0,g}
=
\max_{\{h_{0,i}\}_{i \in S^g}}
\sum_{i \in S^g}
E\!\left\{
u_{1,i}(h_{0,i};h_{0,-g},I)
\right\}
-
\lambda \sum_{i \in S^g} h_{0,i}.
\]

\vspace{0.25cm}
\begin{itemize}
    \item $\lambda>0$ is the per-unit cost of expanding housing capacity.
    \item Each municipality values incumbent period-1 welfare.
    \item Each municipality takes other municipalities' zoning choices as given.
    \item In the decentralized equilibrium, municipalities do not internalize the full regional effect of their zoning choices on housing affordability or transportation returns.
\end{itemize}
\end{frame}

\begin{frame}{Household Period-1 Problem}
A household $\omega$ begins at origin $o$ and chooses residence $i$ and workplace $j$:
\[
u_{1,\omega oij}
=
\max_{c_{ij},h_{ij}}
\frac{B_i}{D_{oij}(I,N)}
\left(\frac{c_{ij}}{\alpha}\right)^{\alpha}
\left(\frac{h_{ij}}{1-\alpha}\right)^{1-\alpha}
\upsilon_{\omega ij}
\]
subject to
\[
c_{ij}+h_{ij}q_i
\leq
w_j-\phi(N_i,h_{0,i})+h_{0,o}q_o .
\]

\vspace{0.2cm}
\begin{itemize}
    \item $B_i$: residential amenity at destination $i$.
    \item $q_i$: housing price at residence location $i$.
    \item $w_j$: wage at workplace location $j$.
    \item $h_{0,o}q_o$: rental income from origin housing owned by the incumbent.
    \item $\phi(N_i,h_{0,i})$: local congestion/crowding cost.
\end{itemize}
\end{frame}

\begin{frame}{Congestion, Moving, and Commuting Frictions}
The spatial friction for a household moving from $o$ to residence $i$ and workplace $j$ is
\[
D_{oij}(I,N)
=
\varepsilon_{oij}
\exp\!\left(
\rho m_{oi}+\kappa \tau_{ij}(I,N)
\right).
\]

\vspace{0.25cm}
\begin{itemize}
    \item $\varepsilon_{oij}$: deterministic bilateral component.
    \item $m_{oi}$: moving cost from origin $o$ to residence $i$.
    \item $\tau_{ij}(I,N)$: commuting cost from residence $i$ to workplace $j$.
    \item Commuting costs decrease with investment and increase with congestion. $\partial\tau/\partial N \geq 0$ (congestion) and $\partial\tau/\partial I \leq 0$ (investment reduces cost).
\end{itemize}

\vspace{0.2cm}


\end{frame}

\begin{frame}{Preference Heterogeneity}
Households draw idiosyncratic preferences over residence-workplace pairs:
\[
\Pr(\upsilon_{\omega ij} \leq z)
=
\exp\!\left(-T_iE_jz^{-\epsilon}\right),
\qquad
z>0,\quad \epsilon>1.
\]

\vspace{0.25cm}
\begin{itemize}
    \item $T_i$ captures residence attractiveness at location $i$.
    \item $E_j$ captures workplace attractiveness at location $j$.
    \item $\epsilon$ governs dispersion in idiosyncratic preferences.
    \item The Fréchet structure gives tractable sorting across residence-workplace pairs.
\end{itemize}

\vspace{0.25cm}
\end{frame}

\begin{frame}{Transportation Agency Problem}
In period 0, after observing the zoning vector $h_0$, the agency chooses network investment:
\[
\max_{\{I_{jk}\}}
\sum_o N_o^0\,\bar U_o(I;h_0)
\]
subject to
\[
\sum_{j,k}\delta_{jk}I_{jk}\leq K,
\qquad
\underline I_{jk}\leq I_{jk}\leq \bar I_{jk}.
\]

\vspace{0.25cm}
Transportation investment enters commuting costs:
\[
\tau_{jk}(I,N)
=
\delta^\tau_{jk}\frac{N_{jk}^{\rho}}{I_{jk}^{\gamma}},
\qquad
\frac{\partial \tau}{\partial N}\geq 0,
\quad
\frac{\partial \tau}{\partial I}\leq 0.
\]

\vspace{0.2cm}
The agency is welfare-oriented, but it treats local zoning capacity as fixed.
\end{frame}

\begin{frame}{Decentralized Equilibrium}
A decentralized equilibrium is a profile
\[
\bigl(
\{h_{0,i}^*\},
I^*,
\{\pi_{oij}^*\},
\{q_i^*,w_j^*\}
\bigr)
\]
such that:
\begin{enumerate}
    \item Households choose residence-workplace pairs optimally.
    \item Housing and labor markets clear at $(q_i^*,w_j^*)$.
    \item Each municipality chooses zoning capacity optimally, taking other municipalities and transportation policy as given.
    \item The transportation agency chooses $I^*$ to maximize regional welfare conditional on $h_0^*$.
\end{enumerate}

\vspace{0.25cm}
\begin{block}{Source of distortion}
The municipality's zoning choice reflects incumbent welfare within its jurisdiction, not the full metropolitan benefits from cheaper housing, better sorting, or stronger transit returns.
\end{block}
\end{frame}

\begin{frame}[shrink=3]{What Local Zoning Omits}
The decentralized municipal choice omits two positive regional externalities:

\vspace{0.12cm}
\begin{block}{1. Cross-municipality welfare externality}
Looser zoning at $i$ lowers local housing prices and makes location $i$ more accessible to households from other municipalities. The local municipality does not fully value these gains.
\end{block}

\vspace{0.12cm}
\begin{block}{2. Transit-zoning externality}
Looser zoning at $i$ raises density near accessible locations. Higher density can raise transit ridership and increase the marginal return to transit investment. A local municipality does not internalize how its zoning affects the regional network choice.
\end{block}

\vspace{0.12cm}
\begin{alertblock}{Main implication}
Decentralized zoning is over-restrictive relative to the coordinated metropolitan optimum:
\[
h_{0,i}^{\mathrm{SP}} > h_{0,i}^{\mathrm{DE}}.
\]
\end{alertblock}
\end{frame}

\begin{frame}{Social Planner Benchmark}
The planner jointly chooses zoning and transportation investment:
\[
\max_{\{h_{0,i}\},\{I_{jk}\}}
\sum_o N_o^0\,\bar U_o(h_0,I)
\]
subject to
\[
\sum_{j,k}\delta_{jk}I_{jk}\leq K,
\qquad
h_{0,i}\geq 0.
\]

\vspace{0.25cm}
\begin{itemize}
    \item The planner internalizes the housing-price effect of zoning across all households.
    \item The planner internalizes how zoning changes optimal transportation investment.
    \item The planner can coordinate density and transit investment in the same locations.
\end{itemize}

\vspace{0.25cm}
\begin{block}{Benchmark comparison}
The welfare gap between the decentralized equilibrium and the planner allocation measures the cost of fragmented land-use control.
\end{block}
\end{frame}

% \begin{frame}{Mechanism: Why Transit Is Especially Affected}
% \begin{columns}[T]
% \column{0.48\textwidth}
% \textbf{Transit}
% \begin{itemize}
%     \item Returns are concentrated around fixed nodes and corridors.
%     \item Ridership depends strongly on nearby residential and workplace density.
%     \item Restrictive zoning near accessible nodes blocks the density response.
%     \item The agency observes weak ridership potential and invests less in transit.
% \end{itemize}

% \column{0.48\textwidth}
% \textbf{Roads}
% \begin{itemize}
%     \item Returns are more spatially diffuse.
%     \item Benefits accrue across many origin-destination pairs.
%     \item Road access can support lower-density, more dispersed locations.
%     \item Road investment is therefore less sensitive to zoning at any one node.
% \end{itemize}
% \end{columns}

% \vspace{0.25cm}
% \[
% \frac{\partial MB^T_{jk}}{\partial h_{0,i^*}}
% >
% \frac{\partial MB^R_{jk}}{\partial h_{0,i^*}}
% \approx 0.
% \]
% \end{frame}

% \begin{frame}[shrink=4]{Putting the Model Together}
% \centering
% \begin{tikzpicture}[
% node distance=0.46cm,
% box/.style={rectangle, rounded corners, draw=blue!65, fill=blue!8, thick,
%             align=center, minimum height=0.54cm, text width=7.6cm, font=\small},
% warn/.style={rectangle, rounded corners, draw=orange!75!black, fill=orange!12, thick,
%              align=center, minimum height=0.54cm, text width=7.6cm, font=\small},
% arrow/.style={-{Latex[length=2mm]}, thick}
% ]
% \node[warn] (zoning) {Municipalities choose over-restrictive zoning};
% \node[box, below=of zoning] (density) {Housing capacity and density near transit-accessible locations are suppressed};
% \node[box, below=of density] (ridership) {Transit ridership and accessibility gains are weaker};
% \node[box, below=of ridership] (returns) {The transportation agency observes lower marginal returns to transit};
% \node[warn, below=of returns] (equilibrium) {Equilibrium tilts toward roads and away from coordinated zoning-plus-transit};

% \draw[arrow] (zoning) -- (density);
% \draw[arrow] (density) -- (ridership);
% \draw[arrow] (ridership) -- (returns);
% \draw[arrow] (returns) -- (equilibrium);
% \end{tikzpicture}

% \vspace{0.16cm}
% \begin{block}{Key takeaway}
% Zoning restrictions do not merely constrain housing supply. They also reshape the perceived return to transportation investment, especially transit investment.
% \end{block}
% \end{frame}

\end{document}
